% =====================================================================
%  Aquisição de sequências no GenBank
%  Prática 3 -- Disciplina Métodos Computacionais (UNESP)
%
%  Compilar com pdfLaTeX (Overleaf: Menu > Compiler > pdfLaTeX).
%
%  VERSÃO: por padrão gera a versão do ALUNO (sem respostas).
%  Para a versão do PROFESSOR (com gabarito), troque \professorfalse
%  por \professortrue na linha indicada abaixo, ou compile com
%     pdflatex "\def\versaoprofessor{}\input{apostila_genbank.tex}"
% =====================================================================
\documentclass[11pt,a4paper]{article}

\newif\ifprofessor
\professorfalse            % <-- troque para \professortrue para incluir o gabarito
\ifdefined\versaoprofessor \professortrue \fi

\usepackage[T1]{fontenc}
\usepackage[utf8]{inputenc}
\usepackage{lmodern}        % Latin Modern: versão vetorial da Computer Modern
\usepackage{textcomp}
\usepackage[provide=*,brazilian]{babel}
\usepackage[a4paper,margin=2.5cm]{geometry}
\usepackage{fancyvrb}
\usepackage{upquote}
\usepackage{array}
\usepackage{booktabs}
\usepackage{longtable}
\usepackage{enumitem}
\usepackage{comment}
\usepackage{url}
\usepackage[hidelinks]{hyperref}

\fvset{xleftmargin=1.5em,fontsize=\small}
\setlist{itemsep=0.2em,topsep=0.4em}
\setlength{\parskip}{0.35em}
\setlength{\emergencystretch}{3em}
\urlstyle{tt}

\newcommand{\arq}[1]{\path{#1}}
\newcommand{\ui}[1]{\textit{#1}}
\newcommand{\tecla}[1]{\textsf{#1}}

\ifprofessor
  \specialcomment{resposta}{\begin{quote}\small\textbf{Gabarito.}\ }{\end{quote}}
\else
  \excludecomment{resposta}
\fi

\title{Aquisição de sequências no GenBank\\[0.3em]
\large Prática 3: BLAST, buscas no site e download de sequências com o R}
\author{Disciplina Métodos Computacionais\\Universidade Estadual Paulista (UNESP)}
\date{\ifprofessor Versão do professor (com gabarito)\else Versão do aluno\fi\\ Novembro de 2026}

\begin{document}
\maketitle

\paragraph{Origem do material.} Esta apostila é uma versão revisada do tutorial
\emph{Aquisição de sequências no GenBank}, publicado em
\url{https://pedrotaucce.github.io/filogenia/}. A sequência das atividades é a mesma; os
códigos em R foram atualizados, e a taxonomia de Hylodidae foi revista para incluir o gênero
\emph{Phantasmarana}. As mudanças estão listadas no Apêndice~\ref{ap:mudancas}.

\paragraph{Arquivos.} A pasta da prática contém:
\begin{itemize}
\item \arq{dados/panthera_onca.fas}: a sequência usada no BLAST (Seção~\ref{sec:blast});
\item \arq{R/}: os scripts \arq{01_preparar_R.R}, \arq{02_baixar_por_nome.R} e
\arq{03_baixar_por_acesso.R}, com todos os comandos desta apostila;
\item \arq{referencia/}: as sequências da Seção~\ref{sec:acesso}, como foram baixadas em 2021,
para usar caso o NCBI esteja inacessível durante a aula.
\end{itemize}

\tableofcontents

% =====================================================================
\section{O que é o GenBank?}
% =====================================================================

O GenBank (\url{https://www.ncbi.nlm.nih.gov/genbank/}), mantido pelo NCBI (\emph{National
Center for Biotechnology Information}, EUA), é a base de dados que reúne as sequências de DNA
disponíveis publicamente. Ele faz parte de um consórcio internacional (INSDC) com o ENA,
europeu, e o DDBJ, japonês: as três bases trocam dados diariamente e compartilham os mesmos
registros e números de acesso.

Quando você elabora um trabalho com sequências de DNA, provavelmente visitará o GenBank para
adquirir sequências antes de começar as análises e, depois, para depositar as sequências que
produziu: o depósito em um repositório público é exigido para a publicação de artigos. Cada
registro recebe um \textbf{número de acesso} (por exemplo, \texttt{KJ961566}) e uma versão
(\texttt{KJ961566.1}), que identificam a sequência de modo permanente.

Uma ferramenta muito útil associada ao GenBank é o \textbf{BLAST} (\emph{Basic Local Alignment
Search Tool}), que compara uma sequência com as sequências da base do NCBI e lista as mais
parecidas. Ele serve, por exemplo, para verificar se uma sequência recém-produzida é de fato do
organismo esperado ou se está contaminada.

% =====================================================================
\section{Usando o BLAST}\label{sec:blast}
% =====================================================================

\begin{enumerate}
\item Abra o arquivo \arq{dados/panthera_onca.fas} em um editor de texto (Bloco de Notas,
Sublime Text, Notepad++). Ele contém uma sequência parcial do gene mitocondrial 16S rRNA
identificada como de onça-pintada (\emph{Panthera onca}) do Pantanal.
\item Acesse o BLAST de nucleotídeos:
\url{https://blast.ncbi.nlm.nih.gov/Blast.cgi?PROGRAM=blastn}.
\item Copie a sequência inteira, incluindo a linha de cabeçalho (\verb|>|), e cole-a no campo
\ui{Enter accession number(s), gi(s), or FASTA sequence(s)}. Outra possibilidade é enviar o
arquivo inteiro pelo botão \ui{Escolher arquivo} (\ui{Or, upload file}).
\item Em \ui{Choose Search Set}, mantenha a base padrão, \ui{Core nucleotide database
(core\_nt)}; em \ui{Program Selection}, mantenha \ui{Highly similar sequences (megablast)}.
\item Clique no botão \ui{BLAST} e aguarde (de alguns segundos a alguns minutos).
\end{enumerate}

Os resultados aparecem na aba \ui{Descriptions}, ordenados da sequência mais parecida para a
menos parecida. As colunas mais importantes são:
\begin{itemize}
\item \ui{Scientific Name}: o organismo da sequência encontrada;
\item \ui{Query Cover}: a porcentagem da sua sequência que foi alinhada;
\item \ui{Per. Ident}: a porcentagem de bases idênticas no trecho alinhado;
\item \ui{E value}: o número de resultados com essa qualidade esperados ao acaso (quanto menor,
melhor; valores como \texttt{0.0} ou \texttt{1e-150} indicam semelhança que não se explica
pelo acaso);
\item \ui{Accession}: o número de acesso da sequência encontrada.
\end{itemize}

\paragraph{Para refletir.} Qual foi a sequência mais próxima da nossa? Se você tivesse acabado
de produzir esta sequência, o que isto poderia significar?

\begin{resposta}
Observe a espécie do primeiro resultado e os valores de \ui{Query Cover} e \ui{Per. Ident}.
Se o primeiro resultado for \emph{Panthera onca} com identidade próxima de 100\%, a
identificação está confirmada. Se for outro organismo (por exemplo, \emph{Homo sapiens}), a
sequência provavelmente resulta de contaminação, e não do animal amostrado; se for outra
espécie do mesmo gênero com identidade alta, pode haver erro de identificação da amostra ou
apenas falta de sequências da espécie na base. O resultado depende do conteúdo atual da base:
rode a busca antes da aula para conferir. Discuta também que o BLAST indica semelhança, não
identidade taxonômica: a identificação das sequências da base depende de quem as depositou.
\end{resposta}

% =====================================================================
\section{Buscando sequências no site do GenBank}
% =====================================================================

Se todas as sequências de DNA públicas estão no GenBank, é lá que procuramos quando precisamos
de alguma. Acesse \url{https://www.ncbi.nlm.nih.gov/nuccore/} e busque pela sequência de algum
organismo e de um fragmento. Se estiver sem ideia, vejamos as sequências do marcador
mitocondrial 16S rRNA da rã-de-unhas-africana (\emph{Xenopus laevis}). Digite no campo de busca:
\begin{Verbatim}
Xenopus laevis 16S
\end{Verbatim}
e clique em \ui{Search}. A página mostra as sequências de \emph{X. laevis} que mencionam o
marcador: são centenas. Repare nos filtros da coluna da esquerda (por exemplo, \ui{Sequence
length}) e no quadro \ui{Search details}, à direita, que mostra como o NCBI interpretou a sua
busca, com \emph{campos} entre colchetes, como \verb|[Organism]|. Os mesmos campos podem ser
usados diretamente:
\begin{Verbatim}
Xenopus laevis[Organism] AND 16S[All Fields] AND 100:5000[SLEN]
\end{Verbatim}
\verb|[SLEN]| é o comprimento da sequência; aqui, de 100 a 5000 pares de bases, o que exclui os
genomas mitocondriais completos (cerca de 17 mil pares de bases).

Baixar centenas de sequências clicando uma a uma daria um trabalho danado! Por isso, foram
criados pacotes para a plataforma R que interagem com o GenBank.

Vamos supor que você encontrou uma espécie nova de rã que parece ser do gênero \emph{Hylodes} e
quer fazer uma árvore com dois marcadores mitocondriais, 16S rRNA e COI (citocromo c oxidase
subunidade I), da família Hylodidae, para posicionar a nova espécie. Baixar no clique todas as
sequências disponíveis para a família seria trabalhoso; vamos fazer isso no R.

% =====================================================================
\section{Um parêntese taxonômico: \emph{Megaelosia} e \emph{Phantasmarana}}\label{sec:taxonomia}
% =====================================================================

Até 2021, a família Hylodidae era dividida em três gêneros: \emph{Crossodactylus},
\emph{Hylodes} e \emph{Megaelosia}. Vittorazzi et al.\ \cite{vittorazzi2021} mostraram que
\emph{Megaelosia} não era monofilético e descreveram o gênero \emph{Phantasmarana} (``rãs
fantasma'', pela raridade com que são encontradas). Desde então, \emph{Megaelosia} contém uma
única espécie, \emph{M. goeldii}, e as demais foram transferidas para \emph{Phantasmarana}, que
hoje tem oito espécies \cite{asw}, duas delas descritas depois da criação do gênero
\cite{desa2022}:

\begin{center}
\begin{tabular}{@{}ll@{}}
\toprule
\emph{Megaelosia} & \emph{Phantasmarana} \\
\midrule
\emph{M. goeldii} (Baumann, 1912) & \emph{P. apuana} (Pombal, Prado, and Canedo, 2003) \\
 & \emph{P. bocainensis} (Giaretta, Bokermann, and Haddad, 1993) \\
 & \emph{P. boticariana} (Giaretta and Aguiar, 1998) \\
 & \emph{P. curucutuensis} de Sá et al., 2022 \\
 & \emph{P. jordanensis} (Heyer, 1983) \\
 & \emph{P. lutzae} (Izecksohn and Gouvêa, 1987) \\
 & \emph{P. massarti} (De Witte, 1930) \\
 & \emph{P. tamuia} de Sá et al., 2022 \\
\bottomrule
\end{tabular}
\end{center}

Isso afeta a busca de sequências de duas maneiras:
\begin{enumerate}
\item A família passa a ter quatro gêneros, e \emph{Phantasmarana} precisa estar na busca.
\item As sequências depositadas antes de 2021 foram descritas como \emph{Megaelosia}. A
Taxonomia do NCBI atualiza o campo \texttt{ORGANISM} dos registros, que é o usado nas buscas
com \verb|[Organism]|; mas o texto descritivo do registro (a linha que aparece no cabeçalho do
FASTA) foi escrito por quem depositou a sequência e pode continuar com o nome antigo. Ao montar
uma matriz, confira sempre os nomes.
\end{enumerate}

% =====================================================================
\section{Usando o R para adquirir sequências}
% =====================================================================

\subsection{Preparando o R}

Usaremos o R (\url{https://www.r-project.org/}) por meio do RStudio. Recomenda-se criar um
\emph{projeto} do RStudio para a prática (\ui{File > New Project > New Directory}): o RStudio
passa a usar a pasta do projeto como pasta de trabalho, e todos os arquivos são gravados ali,
sem a necessidade de \verb|setwd()| com caminhos que só existem no seu computador. Copie para a
pasta do projeto as pastas \arq{R}, \arq{dados} e \arq{referencia}.

Precisamos de dois pacotes: \texttt{ape} \cite{paradis2019}, para ler e gravar sequências, e
\texttt{rentrez} \cite{winter2017}, para conversar com as bases do NCBI. Abra o script
\arq{R/01_preparar_R.R} e execute-o (no RStudio, \tecla{Ctrl+Enter} executa a linha atual;
\tecla{Ctrl+Shift+Enter}, o script inteiro):
\begin{Verbatim}
install.packages(c("ape", "rentrez"))
library(ape)
library(rentrez)
packageVersion("ape")
packageVersion("rentrez")
\end{Verbatim}
\verb|library()| carrega o pacote; se ele não estiver instalado ou não funcionar, aparece uma
mensagem de erro (por exemplo, \texttt{there is no package called 'ape'}). Se nenhum erro
aparecer e as duas últimas linhas mostrarem as versões instaladas (algo como
\verb|[1] '5.8.1'|), você pode seguir adiante. Anote as versões: elas fazem parte da
documentação de um trabalho. A instalação só precisa ser feita uma vez em cada computador; nas
sessões seguintes basta carregar os pacotes com \verb|library()|.

\paragraph{Limite de acessos do NCBI.} O NCBI aceita até 3 requisições por segundo de cada
endereço de internet. Em uma sala em que todos usam a mesma rede, esse limite pode ser atingido,
e aparecerá um erro como \texttt{HTTP 429 Too Many Requests}: espere alguns segundos e execute
o comando de novo. Quem cria uma conta no NCBI pode gerar uma chave pessoal da API (\emph{API
key}), que eleva o limite para 10 requisições por segundo, e informá-la ao R com
\verb|set_entrez_key("...")|. A chave é pessoal: não a publique junto com scripts.

\subsection{Baixando as sequências por nome}

Abra o script \arq{R/02_baixar_por_nome.R}. Vamos executá-lo por partes, entendendo cada uma.

\paragraph{Parte 1: o que buscar.}
\begin{Verbatim}[fontsize=\footnotesize]
library(ape)
library(rentrez)
dir.create("sequencias", showWarnings = FALSE)

generos <- c("Crossodactylus", "Hylodes", "Megaelosia", "Phantasmarana")

termo_coi <- paste("(COI[All Fields] OR COX1[All Fields]",
                   'OR "cytochrome oxidase subunit 1"[All Fields]',
                   'OR "cytochrome oxidase subunit I"[All Fields]',
                   'OR "cytochrome c oxidase subunit I"[All Fields])')
marcadores <- c("16S" = "16S[All Fields]", "COI" = termo_coi)
filtro_tamanho <- "100:5000[SLEN]"
termo_generos <- paste0("(", paste0(generos, "[Organism]", collapse = " OR "), ")")
\end{Verbatim}
O COI aparece no GenBank com vários nomes; por isso o termo junta as variantes com \verb|OR|
(\verb|paste()| une os pedaços de texto, separados por espaço). A última linha monta o texto
\verb|(Crossodactylus[Organism] OR Hylodes[Organism] OR ...)|: digite \verb|termo_generos| no
console para vê-lo.

\paragraph{Parte 2: uma função para buscar e baixar.}
\begin{Verbatim}[fontsize=\footnotesize]
baixar_fasta <- function(termo, arquivo, lote = 200) {
  busca <- entrez_search(db = "nuccore", term = termo, use_history = TRUE)
  if (busca$count == 0) return(0)
  if (file.exists(arquivo)) file.remove(arquivo)
  for (inicio in seq(0, busca$count - 1, by = lote)) {
    fasta <- entrez_fetch(db = "nuccore", web_history = busca$web_history,
                          rettype = "fasta", retstart = inicio, retmax = lote)
    cat(fasta, file = arquivo, append = TRUE)
    Sys.sleep(0.5)
  }
  busca$count
}
\end{Verbatim}
\verb|entrez_search()| faz a busca na base \texttt{nuccore} (nucleotídeos) e, com
\verb|use_history = TRUE|, deixa a lista de resultados guardada no servidor do NCBI.
\verb|entrez_fetch()| baixa as sequências em formato FASTA, em lotes de 200 (\verb|retstart|
indica a partir de qual resultado), o que funciona mesmo quando há milhares de resultados. O
arquivo antigo é apagado antes, para que executar o script de novo não duplique as sequências,
e \verb|Sys.sleep(0.5)| faz uma pausa entre os lotes, respeitando o limite do NCBI.

\paragraph{Parte 3: baixar um arquivo por marcador.}
\begin{Verbatim}[fontsize=\footnotesize]
for (m in names(marcadores)) {
  termo <- paste(termo_generos, "AND", marcadores[[m]], "AND", filtro_tamanho)
  arquivo <- file.path("sequencias", paste0("hylodidae_", m, ".fasta"))
  n <- baixar_fasta(termo, arquivo)
  message(m, ": ", n, " sequências gravadas em ", arquivo)
}
\end{Verbatim}
Ao final, a pasta \arq{sequencias} contém \arq{hylodidae_16S.fasta} e
\arq{hylodidae_COI.fasta}, cada um com todas as sequências de Hylodidae do marcador. Anote
quantas sequências foram gravadas: o número cresce à medida que novas sequências são
depositadas.

\paragraph{Parte 4: quantas sequências há de cada gênero?} O script repete a busca para cada
gênero e marcador, pedindo só a contagem (\verb|retmax = 0|), e monta uma tabela:
\begin{Verbatim}[fontsize=\footnotesize]
contagem <- sapply(generos, function(g) {
  sapply(names(marcadores), function(m) {
    Sys.sleep(0.5)
    termo <- paste(paste0(g, "[Organism]"), "AND", marcadores[[m]], "AND", filtro_tamanho)
    entrez_search(db = "nuccore", term = termo, retmax = 0)$count
  })
})
contagem
\end{Verbatim}

\paragraph{Parte 5: conferir os arquivos.}
\begin{Verbatim}
s16 <- read.FASTA("sequencias/hylodidae_16S.fasta")
coi <- read.FASTA("sequencias/hylodidae_COI.fasta")
length(s16)
head(names(s16))
summary(lengths(s16))
\end{Verbatim}
\verb|length()| informa o número de sequências; \verb|names()|, os cabeçalhos; e
\verb|summary(lengths())|, o resumo dos comprimentos em pares de bases.

\paragraph{Parte 6: nomes curtos.} Os cabeçalhos do GenBank são longos
(\texttt{KJ961566.1 Megaelosia boticariana voucher CFBHT00425 cytochrome oxidase\ldots}).
Para transformá-los em \verb|Genero_especie_acesso|, usamos a mesma expressão regular da
Prática~1 (no R, cada \verb|\| é escrita \verb|\\|, e \verb|perl = TRUE| permite o grupo de não
captura \verb|(?:...)|, que aqui ignora o aviso \texttt{UNVERIFIED:} presente em alguns
registros):
\begin{Verbatim}[fontsize=\footnotesize]
renomear <- function(nomes) {
  sub("^([A-Z]+_?[0-9]+)\\.[0-9]+ (?:UNVERIFIED: )?([A-Z][a-z]+) ([a-z]+).*$",
      "\\2_\\3_\\1", nomes, perl = TRUE)
}
names(s16) <- renomear(names(s16))
names(coi) <- renomear(names(coi))
names(s16)[duplicated(names(s16))]
write.FASTA(s16, "sequencias/hylodidae_16S_nomes.fasta")
write.FASTA(coi, "sequencias/hylodidae_COI_nomes.fasta")
\end{Verbatim}
A penúltima linha verifica se há nomes repetidos (deve mostrar \verb|character(0)|).

\subsection{Baixando as sequências por número de acesso}\label{sec:acesso}

Suponhamos que acaba de sair a descrição de uma nova espécie de \emph{Phantasmarana}, como
\emph{P. curucutuensis} e \emph{P. tamuia}, descritas em 2022 \cite{desa2022}, e que os autores
fizeram uma ótima árvore, à qual você gostaria de acrescentar a sua espécie nova de
\emph{Hylodes}. Infelizmente, os autores não disponibilizaram a matriz usada no artigo, mas o
material suplementar traz a lista de números de acesso de todas as sequências. Nesse caso, você
pode baixar exatamente as mesmas sequências. Abra o script \arq{R/03_baixar_por_acesso.R}:
\begin{Verbatim}[fontsize=\footnotesize]
library(ape)
ids_16S <- c("KY007117.1", "KU215902.1", "KM390791.1", "AY843579.1", "AY263235.1",
             "MF624225.1", "KY627916.1", "KJ961575.1", "KJ961577.1", "MF624238.1",
             "AY843689.1", "KJ961586.1", "KJ961578.1", "DQ283072.1", "KJ961583.1",
             "KM390793.1", "KM390794.1", "KU495250.1", "KY002953.1", "KY627911.1",
             "KY627913.1", "KY627914.1")
ids_COI <- c("KU494387.1", "DQ502738.1", "KU494457.1", "KC593345.1", "KJ961562.1",
             "KJ961561.1", "KJ961558.1", "KJ961557.1", "KJ961556.1", "DQ502873.1",
             "KJ961565.1", "KU494594.1", "KJ961566.1", "DQ502839.1")
s16 <- read.GenBank(ids_16S)
coi <- read.GenBank(ids_COI)
\end{Verbatim}
A função \verb|read.GenBank()|, do \texttt{ape}, baixa as sequências e guarda, no atributo
\texttt{species}, o nome do organismo segundo a taxonomia \emph{atual} do NCBI. Compare esse
nome com a descrição original de cada registro:
\begin{Verbatim}[fontsize=\footnotesize]
data.frame(acesso = names(coi),
           organismo = attr(coi, "species"),
           descricao = substr(attr(coi, "description"), 1, 60))
\end{Verbatim}
Por fim, damos às sequências nomes no formato \verb|Genero_especie_acesso|, usando o nome atual
do organismo, e gravamos os arquivos:
\begin{Verbatim}[fontsize=\footnotesize]
names(s16) <- paste(attr(s16, "species"), sub("\\.[0-9]+$", "", names(s16)), sep = "_")
names(coi) <- paste(attr(coi, "species"), sub("\\.[0-9]+$", "", names(coi)), sep = "_")
write.FASTA(s16, "sequencias/16S_ids.fasta")
write.FASTA(coi, "sequencias/COI_ids.fasta")
\end{Verbatim}
Agora você tem no seu computador os arquivos \arq{16S_ids.fasta} e \arq{COI_ids.fasta}, com
as sequências desejadas. O script traz, comentada, a versão com o \texttt{rentrez}
(\verb|entrez_fetch()|), que grava o FASTA exatamente como o NCBI o fornece.

\paragraph{Sem acesso ao NCBI?} Se a rede da sala bloquear o NCBI ou o serviço estiver fora do
ar, use os arquivos \arq{referencia/16S_ids_download_2021.fasta} e
\arq{referencia/COI_ids_download_2021.fasta}: são as mesmas sequências, baixadas em 2021, com os
cabeçalhos da época.

Com as sequências em mãos, o próximo passo é o alinhamento.

% =====================================================================
\section*{Para refletir}
\addcontentsline{toc}{section}{Para refletir}
% =====================================================================

\begin{enumerate}
\item Na tabela \verb|contagem|, quantas sequências de 16S e de COI há para \emph{Megaelosia}
e para \emph{Phantasmarana}? Que espécie(s) aparece(m) quando se busca
\verb|Megaelosia[Organism]|?
\item Procure, nos cabeçalhos de \arq{hylodidae_COI.fasta} (ou nos nomes gerados pela
função \verb|renomear()|), sequências de \emph{Phantasmarana}. Elas aparecem com esse nome?
O que isso indica sobre a confiabilidade das linhas de descrição dos registros?
\item Na Seção~\ref{sec:acesso}, compare as colunas \texttt{organismo} e \texttt{descricao}
da tabela. Quais registros mudaram de nome?
\item A lista de números de acesso inclui \emph{Limnomedusa macroglossa}, que não pertence a
Hylodidae. Por que ela foi incluída?
\item Por que excluímos sequências com mais de 5000 pares de bases? O que aconteceria se um
genoma mitocondrial completo entrasse no arquivo do 16S?
\item Executar o script \arq{02_baixar_por_nome.R} hoje e daqui a um ano produziria o mesmo
arquivo? Como registrar, em um trabalho, quais sequências foram usadas?
\end{enumerate}

\begin{resposta}
1. Os números mudam com o tempo; espera-se que \emph{Megaelosia} tenha apenas sequências de
\emph{M. goeldii}, e \emph{Phantasmarana}, as das demais espécies (inclusive as depositadas
como \emph{Megaelosia}), pois a busca por \verb|[Organism]| usa a taxonomia atual.
2. Sequências depositadas antes de 2021 tendem a manter ``Megaelosia'' na descrição; os nomes
gerados a partir dela (por exemplo, \texttt{Megaelosia\_boticariana\_KJ961566}) ficam
desatualizados. A descrição é texto livre do depositante; o campo \texttt{ORGANISM} é o que
acompanha a taxonomia. Por isso a Seção~\ref{sec:acesso} usa \verb|attr(x, "species")|.
3. Esperado (se a Taxonomia do NCBI estiver atualizada): \texttt{MF624238} (\emph{P.
jordanensis}), \texttt{KJ961586} e \texttt{KJ961566} (\emph{P. boticariana}) e
\texttt{KU494594} (\emph{P. apuana}); \texttt{DQ283072} e \texttt{DQ502839} continuam
\emph{Megaelosia goeldii}. O script \arq{R/verificar_pratica3.R} confere isso.
4. É o grupo externo (Alsodidae, família proximamente relacionada a Hylodidae), usado para
enraizar a árvore.
5. Um genoma mitocondrial tem cerca de 17 mil pares de bases e contém todos os genes; no
alinhamento do 16S, ele introduziria milhares de posições sem homologia com as demais
sequências. Ele pode ser aproveitado, mas extraindo-se apenas o trecho do 16S.
6. Não: novas sequências são depositadas continuamente, e registros podem ser atualizados. Por
isso, publica-se a lista de números de acesso (com a versão) de todas as sequências usadas,
como na Seção~\ref{sec:acesso}, além da data da busca e do termo usado.
\end{resposta}

% =====================================================================
\clearpage
\appendix
\section{Problemas comuns}\label{ap:problemas}
% =====================================================================

\begin{description}[style=nextline,leftmargin=1.5em]
\item[\texttt{HTTP 429} ou \texttt{Too Many Requests}] Limite de requisições do NCBI atingido
(comum quando a turma toda usa a mesma rede). Aguarde alguns segundos e repita; se possível, use
uma chave da API (\verb|set_entrez_key()|).
\item[\texttt{HTTP 500}, \texttt{HTTP 502} ou tempo esgotado] Instabilidade do NCBI. Repita
mais tarde ou use os arquivos da pasta \arq{referencia}.
\item[\texttt{cannot open file}] O arquivo não está na pasta de trabalho. Confira com
\verb|getwd()| e \verb|list.files()|; no RStudio, abra o projeto da prática.
\item[\texttt{there is no package called}] O pacote não foi instalado: execute
\arq{01_preparar_R.R}.
\item[Arquivo com sequências repetidas] Ocorria na versão anterior do tutorial, que acrescentava
as sequências ao arquivo existente a cada execução. Os scripts atuais apagam o arquivo antes de
gravar.
\end{description}

% =====================================================================
\section{Alterações em relação à versão anterior}\label{ap:mudancas}
% =====================================================================

\begin{enumerate}
\item \textbf{Taxonomia.} Inclusão de \emph{Phantasmarana} \cite{vittorazzi2021,desa2022} na
busca por nome e no cenário da busca por número de acesso; discussão sobre nomes antigos nas
descrições dos registros.
\item \textbf{Arquivos de saída.} O código anterior usava \verb|paste(i, "hylodidae_16S.fasta",|
\verb|sep = "_")| dentro do laço, o que gravava um arquivo por gênero
(\arq{Hylodes_hylodidae_16S.fasta} etc.), e não os dois arquivos mencionados no texto. Agora a
busca reúne os gêneros e grava um arquivo por marcador.
\item \textbf{Sequências duplicadas.} Com \verb|append = TRUE|, executar o script de novo
duplicava as sequências. Agora o arquivo é recriado.
\item \textbf{Termos de busca.} \verb|16S [ALL]| e \verb|coi [ALL]| foram substituídos por
termos com os nomes alternativos do COI e pelo filtro de comprimento \verb|[SLEN]|, que exclui
genomas mitocondriais completos.
\item \textbf{Downloads grandes.} Uso do \emph{web history} do NCBI e de lotes de 200
sequências, com pausa entre os lotes, no lugar de um único pedido com até 10~000 identificadores.
\item \textbf{Pasta de trabalho.} \verb|setwd("C:/caminho/...")| foi substituído por um projeto
do RStudio e caminhos relativos.
\item \textbf{Instalação.} \verb|install.packages()| com os dois pacotes de uma vez e conferência
com \verb|library()| e \verb|packageVersion()|, no lugar de \verb|installed.packages()|, que só
verifica se a pasta do pacote existe, e não se ele funciona.
\item \textbf{Nomes.} Renomeação para \verb|Genero_especie_acesso| com expressão regular (ligação
com a Prática 1) e, na busca por acesso, com o nome atual do organismo
(\verb|read.GenBank()|).
\item \textbf{BLAST.} Instruções atualizadas para a página atual, que usa por padrão a base
\emph{core\_nt} desde 2024, e descrição das colunas do resultado; a captura de tela foi
substituída pela descrição em texto.
\item \textbf{Arquivos de referência.} Inclusão das sequências baixadas em 2021, para uso sem
acesso ao NCBI, e de um script de verificação para o professor.
\end{enumerate}

\begin{thebibliography}{9}
\bibitem{asw} Frost, D.~R. 2026. \emph{Amphibian Species of the World: an Online Reference},
versão 6.2. American Museum of Natural History, Nova York.
\url{https://amphibiansoftheworld.amnh.org/}. Acesso em 9 nov. 2026.
\bibitem{desa2022} de Sá, F.~P., Condez, T.~H., Lyra, M.~L., Haddad, C.~F.~B. \& Malagoli, L.~R.
2022. Unveiling the diversity of giant neotropical torrent frogs (Hylodidae): phylogenetic
relationships, morphology, and the description of two new species. \emph{Systematics and
Biodiversity} 20(1): 2039318.
\bibitem{paradis2019} Paradis, E. \& Schliep, K. 2019. ape 5.0: an environment for modern
phylogenetics and evolutionary analyses in R. \emph{Bioinformatics} 35: 526--528.
\bibitem{vittorazzi2021} Vittorazzi, S.~E., Augusto-Alves, G., Neves-da-Silva, D.,
Carvalho-e-Silva, A.~M.~P.~T., Recco-Pimentel, S.~M., Toledo, L.~F., Lourenço, L.~B. \&
Bruschi, D.~P. 2021. Paraphyly in the giant torrent-frogs (Anura: Hylodidae:
\emph{Megaelosia}) and the description of a new genus. \emph{Salamandra} 57(2): 274--284.
\bibitem{winter2017} Winter, D.~J. 2017. rentrez: an R package for the NCBI eUtils API.
\emph{The R Journal} 9(2): 520--526.
\end{thebibliography}

\end{document}
